\subsection{III)}
{\noindent\large\begin{align*}\displaystyle F(x, y, z) &= \sum(0, 2, 4, 6)\end{align*}}
\subsubsection{Obtener la mínima expresión booleana}
\begin{center}
\begin{tabular}{c|ccc|c}
\rowcolor[HTML]{FFCCC9} 
\cellcolor[HTML]{9AFF99}$m$ & $x$ & $y$ & $z$ & \cellcolor[HTML]{FFFFC7}$F(x,y,z)$ \\ \hline
0                           & 0   & 0   & 0   & 1                                  \\
1                           & 0   & 0   & 1   &                                    \\
2                           & 0   & 1   & 0   & 1                                  \\
3                           & 0   & 1   & 1   &                                    \\
4                           & 1   & 0   & 0   & 1                                  \\
5                           & 1   & 0   & 1   &                                    \\
6                           & 1   & 1   & 0   & 1                                  \\
7                           & 1   & 1   & 1   &                                   
\end{tabular}
\end{center}

\begin{center}
    \begin{karnaugh-map}[2][4][1][$z$][$x\qquad y$]
        \minterms{0,2,4,6}
        \implicant{0}{4}
    \end{karnaugh-map}
\end{center}
\begin{itemize}\item En el bloque permanece constante $z$, por lo que se deja y hay que negarlo.\item La función simplificada queda:\end{itemize}
{\large\begin{align*}
    F(x,y,z) &= \overline{z}
\end{align*}}
\imagen[]{10cm}{./imagenes/ejercicio3Falstad.png}
\subsubsection{Descripción en Verilog}
\begin{lstlisting}
module descripcionVerilogEjercicio3(
    input inputX,
    input inputY,
    input inputZ,
    output outputFunction
    );
	 assign outputFunction = ~inputZ;

endmodule
\end{lstlisting}
\subsubsection{RTL output}
\imagen[Bloque RTL]{15cm}{./imagenes/rtlEjercicio3Verilog.png}
\imagen[Implementación con compuertas]{15cm}{./imagenes/rtlCompuertasEjercicio3Verilog.png}

\subsubsection{Simulación temporal}
\noindent Código del módulo Verilog Test Fixture:
\begin{lstlisting}
module verilogTestFixtureEjercicio3;

	// Inputs
	reg inputX;
	reg inputY;
	reg inputZ;

	// Outputs
	wire outputFunction;

	// Instantiate the Unit Under Test (UUT)
	descripcionVerilogEjercicio3 uut (
		.inputX(inputX), 
		.inputY(inputY), 
		.inputZ(inputZ), 
		.outputFunction(outputFunction)
	);

	initial begin
		// Initialize Inputs
		inputX = 0;
		inputY = 0;
		inputZ = 0;

		// Wait 100 ns for global reset to finish
		#100;
      inputX = 0; inputY = 0; inputZ = 1;
		#100;
      inputX = 0; inputY = 1; inputZ = 0;
		#100;
      inputX = 0; inputY = 1; inputZ = 1;
		#100;
      inputX = 1; inputY = 0; inputZ = 0;
		#100;
      inputX = 1; inputY = 0; inputZ = 1;
		#100;
      inputX = 1; inputY = 1; inputZ = 0;
		#100;
      inputX = 1; inputY = 1; inputZ = 1;		
		// Add stimulus here

	end
      
endmodule
\end{lstlisting}
\imagen[Simulación temporal]{15cm}{./imagenes/simulacionTemporalEjercicio3.png}
